\par\bigskip\noindent\textbf{K101 supporting proof: The actual old cap, graph roots and full annular groupoid}\par\nobreak
Working calculation owned by blind\_control, complementary to
\path{part3_e6_selected_infinity_working.md} (the source/Miller/collar calculation).
The maps below have been derived on their displayed charts. The complementary
root corner, the full Cartier graph there, and its completed suspension
attachment are still being calculated. The actual ordinary quotient,
point-line chain and cap-domain overlaps are written below. No return
or compact global boundary is assigned by this note.

\par

\medskip\noindent\textbf{EIR.1. Both actual native roots and their directed source maps}

\par\nobreak

Put
\[
 e=1/w,\quad d=1-3e^2,\quad A=3e^2-2,\quad
 C=(1-e^2)(1+3e^2),\quad C_1=4d^2-C.
\]
The actual source equation, with the same marked section, is
\[
 y^2=x^3+Ax^2+dtx+e^2t^2,\qquad G=(0,et),\qquad
 q_e=t-(1-e^2).
 \tag{EIR.1}
\]
Its original coordinates are
\[
 x_o=e^2d^{-2}x,\quad y_o=e^3d^{-3}y,\quad
 t_o=e^4d^{-3}t,\quad r=e^2(3e^2-2)/d^2.
 \tag{EIR.2}
\]
In particular the original source parameter has local order two at this
ordered point. The finite selected height and the height at infinity are
different coordinates and have the exact relations
\[
 q_o=e^4d^{-3}q_e,\qquad q_{\infty,o}=e^{-4}d^3q_{\infty,e}.
 \tag{EIR.3}
\]

Use the following uniform root notation:
\[
 \begin{array}{c|cccc}
 (a,b)&m&k&r_*&c\\\hline
 (2,3)&3l&2l&l&0\quad(l\ge1)\\
 (6,1)&3l+2&2l+1&l+1&1\quad(l\ge0).
 \end{array}
\]
Thus \(m=3r_*-c\), \(k=2r_*-c\). Set
\(q_{\infty,e}=z^m\), \(x_0=z^{2k}X\), \(y_0=z^{3k}Y\).
The full native root, before further modification at the joint origin, is
\[
 Y^2=X^3+A z^{2r_*}X^2+d z^{r_*+c}X+e^2z^{2c},
 \qquad G'=(0,ez^c).
 \tag{EIR.4}
\]
Let \(J\) denote the native-to-original map. Choose the analytic root of
the nonzero unit \(d\), and choose the power of \(e\) on its logarithmic
parameter cover. Its complete coordinates are
\[
 \begin{array}{c|ccc}
 &z_o&X_o&Y_o\\\hline
 a=2&e^{-4/m}d^{3/m}z&e^{-2/3}X&e^{-1}Y\\
 a=6&e^{-4/m}d^{3/m}z&
 e^{-2(l+2)/m}d^{2/m}X&e^{-3(l+2)/m}d^{3/m}Y .
 \end{array}
 \tag{EIR.5}
\]
Substitution into the original K96 equations gives (EIR.4) exactly. In
both rows it sends the actual marked section to the original marked
section and commutes with the finite root deck. If
\(u=\log e/(2\pi i)\), continuation once gives
\[
 J_{u+1}=\mathfrak D_o^{-4}J_u,
 \qquad J_{u+1}^{-1}=\mathfrak D_e^{4}J_u^{-1}.
 \tag{EIR.6}
\]
The first formula is the new-to-old arrow. The second is its inverse.
For the smallest reduced cap \(m=2\), \(\mathfrak D^{-4}=1\), so
this actual coordinate comparison itself is already closed on the
parameter puncture. In that case the later compensator is the explicit
deck \(C=A^{-4}=B^{-2n}\); no logarithm is required to close this
particular root-coordinate arrow.

The same source map carries the canonical point-line section literally.
Writing \(\alpha=e/d\), the actual Miller functions satisfy
\(f_{a,o}=\alpha^{a-1}f_{a,e}\). Hence the two frame dictionaries are
\[
 V_{{\rm pt},o}=V_{{\rm pt},e},\qquad
 V_{{\rm tensor},o}=\alpha^{1-a}V_{{\rm tensor},e},
\]
\[
 q_o^h Z_o^v V_{{\rm tensor},o}^j
 =e^{4h+(1-a)j}d^{-3h+(a-1)j}
      q_e^h Z_e^v V_{{\rm tensor},e}^j.
 \tag{EIR.7}
\]
All these powers are part of the actual chart map. They are not
discarded as extending units at \(e=0\).

The full source calculation gives
\[
 B_{{\rm sel},o}=e^4\widehat B_e,\qquad
 \widehat B_e=e^{-6}(q_e/t^3)B_{a,b,e}.
 \tag{EIR.8}
\]
Consequently the same original fixed scalar has native value
\(\lambda_e=e^4\lambda\). The complete root/line map on the punctured
parameter disc carries this equality, with its actual Cartier factors.
The total finite-cap generator is carried; a separate choice of
\(\lambda_e^{1/m}\) and a separate root of \(\widehat B_e\) each has
parameter monodromy, which cancels in that full generator.

\par

\medskip\noindent\textbf{EIR.2. Exact joint formula for the formal coordinate of three times G}

\par\nobreak

In (EIR.4), put
\[
 g=ez^c,\qquad d_G=\frac d{2e}z^{r_*},\qquad
 s_G=\frac C{4e^2}z^{2r_*}.
\]
The literal tangent and next chord give
\[
 x(2G')=s_G,\quad y(2G')=-g-d_Gs_G,
\]
\[
 x(3G')=\frac{4d_Gg}{s_G}+\frac{4g^2}{s_G^2},\qquad
 y(3G')=-g+\left(d_G+\frac{2g}{s_G}\right)x(3G').
 \tag{EIR.9}
\]
In the actual identity formal coordinate \(\tau=-X/Y\), these simplify
to the following single rational function for both progressions:
\[
 \tau(3G')=
 -\frac{8eCz^k\bigl(dCz^m+8e^4\bigr)}
 {C^2C_1z^{2m}+96dCe^4z^m+512e^8}.
 \tag{EIR.10}
\]
In particular, at fixed \(e\ne0\), its leading coefficient is
\[
 \tau(3G')=-\frac C{8e^3}z^k+\text{higher powers of }z.
 \tag{EIR.11}
\]
This is a literal coefficient pole. An analytic formal-coordinate
function on a fixed joint bidisc would have holomorphic Taylor
coefficients in \(e\), so (EIR.11) precludes that particular
extension before modifying the joint chart.

For the good row \(R=[m]^{-1}(3G')\). For the reduced row the
actual division is
\(R=2G'-[m]^{-1}((2m-1)G')\); its central component is two in the
order-three component group. Its identity-component multiple is
\[
 3R=[m]^{-1}(3G').
 \tag{EIR.12}
\]
Indeed both sides lie in the identity neighbourhood and their
\(m\)-multiples are the same section, so the chosen local inverse
of \([m]\) gives this exact equality. It follows that the good
\(R\) has leading coordinate \(-Cz^{2l}/(8me^3)\), while the
reduced identity-component section \((m+4)R\) has leading coordinate
\[
 -\frac{(l+2)C}{8me^3}z^{2l+1},\qquad m+4=3(l+2).
 \tag{EIR.13}
\]
Multiplying the component by three has removed its component class;
the displayed analytic pole remains and is carried into the joint graph.

\par

\medskip\noindent\textbf{EIR.3. The actual joint base graph and the old-root chart}

\par\nobreak

The root-height equality (EIR.3) is
\[
 z_o^m=d^3z^m/e^4.
 \tag{EIR.14}
\]
The common base graph is the normalization of the blowup ideal
\((e^4,z^m)\). In valuation coordinates \((v_e,v_z)\), its
distinguished ray is \((m,4)\). Its primitive generator is
\((m/g_0,4/g_0)\), where \(g_0=\gcd(m,4)\).

Retain the full cover
\[
 e=\eta^m,\qquad z=\eta^4\zeta.
 \tag{EIR.15}
\]
Then \(z_o=d^{3/m}\zeta\). The actual source coordinates on this
chart are
\[
 \begin{array}{c|cc}
 &X&Y\\\hline
 a=2&\eta^{2l}U&\eta^{3l}V\\
 a=6&d^{-2/m}\eta^{2(l+2)}U&
        d^{-3/m}\eta^{3(l+2)}V .
 \end{array}
 \tag{EIR.16}
\]
Here \((U,V,z_o)\) are precisely the original K96 root coordinates;
its parameter is \(r=e^2A/d^2\). Thus this one chart carries the
original root at \(r=0\) with every original marked section. It
does so by the literal source equation, rather than by matching
the central Kodaira symbol.

The full \(\mu_m\) cover action is
\[
 (\eta,\zeta,U,V)\longmapsto
 (\zeta_m\eta,\zeta_m^{-4}\zeta,
   \zeta_m^{-2b_*}U,\zeta_m^{-3b_*}V),
 \quad b_*=l\ (a=2),\quad b_*=l+2\ (a=6).
 \tag{EIR.17}
\]
On the original root this is exactly \(\mathfrak D_o^{-4}\).
The congruences follow directly from \(m=3l\) or \(m=3l+2\).
The subgroup of order \(g_0\) acts trivially on
\((\zeta,U,V)\) and acts on \(\eta\); the primitive weighted
chart retains the quotient of this cover with that subgroup.
This accounts for the finite kernel rather than omitting it.

Formula (EIR.10) also reads, with \(Q=z^m/e^4\),
\[
 \tau(3G')=-8C\frac{z^k}{e^3}
       \frac{dCQ+8}{C^2C_1Q^2+96dCQ+512}.
 \tag{EIR.18}
\]
The last quotient is a nonzero analytic unit near \(Q=0,e=0\).
The exact coordinate comparison gives
\[
 \tau_o=e^{1/3}\tau_e\quad(a=2),\qquad
 \tau_o=e^{(l+2)/m}d^{-1/m}\tau_e\quad(a=6).
 \tag{EIR.19}
\]
Combining (EIR.18)-\/-(EIR.19) recovers the regular original
formal coefficient \(-Cz_o^k/(8d^2)\). Conversely, the domain
\(|z_o|<\delta\) becomes
\[
 |z|<\delta |e|^{4/m}|d|^{-3/m}.
 \tag{EIR.20}
\]
This is the precise domain on which the original formal logarithm
has so far been transported. It does not supply a fixed native
bidisc at the joint origin.

\par

\medskip\noindent\textbf{EIR.4. Cartier factors on the graph and the remaining root corner}

\par\nobreak

For the good root put \(u=z^l\). On the punctured parameter chart
the marked second point and the two old Cartier coordinates are
\[
 x(2G')=\frac C{4e^2}u^2,\quad
 w_{P,e}=X-\frac C{4e^2}u^2,\quad w_{O,e}=-X/Y,
\]
\[
 u_o=e^{-4/3}d\,u,\quad
 w_{P,o}=e^{-2/3}w_{P,e},\quad
 w_{O,o}=e^{1/3}w_{O,e}.
 \tag{EIR.21}
\]
The original coefficient is \(s_o=C/(4d^2)\), so substitution
proves the first line-coordinate equality including the actual
second section. In the K96 exceptional corrections, a generator
\(w_P^r u^s\) consequently carries
\(e^{-(2r+4s)/3}d^s\), whereas \(w_O^r u^s\) carries
\(e^{(r-4s)/3}d^s\). These are the full monomial factors in
\(\beta^*N_2=\pi^*M_2+E_P-E_O\).

The old-root chart (EIR.16) has the additional local O-arrow
\[
 \tau_e=\eta^{-l}\tau_o\quad(a=2),\qquad
 \tau_e=d^{1/m}\eta^{-(l+2)}\tau_o\quad(a=6).
 \tag{EIR.22}
\]
Its graph at the old zero section is the graph of the ideal
\((\tau_o,\eta^{b_*})\), with the displayed unit in the reduced
row. For the reduced line the two original marked equations are
\(\tau_o\) and \(\tau_o-\tau_o(6G')\); their actual ratio is
retained on this graph. The joint intersections with
\(\zeta=0\), and the complementary \(\zeta=\infty\) corner,
require their full resolved Cartier comparison. No equal-
coefficient claim at those intersections is made here.

\par

\medskip\noindent\textbf{EIR.5. Full cap endpoint, the actual suspension reindex, and its domain}

\par\nobreak

On the original root let
\[
 A=R_\lambda^n\zeta_m^d\widetilde F^{\,n}\mathfrak D,
 \quad B=\lambda B_0,\quad
 \widetilde F^m=B_0,\quad A^m=B^n,
 \quad n=\pm1.
 \tag{EIR.23}
\]
The symbols here are the actual original K96 source/Cartier lifts;
they are holomorphic over the original \(r=0\) disc. The base-
preserving word for the directed endpoint (EIR.6) is
\[
 C=A^{-4}\mathfrak D^4
   =R_\lambda^{-4n}\zeta_m^{-4d}\widetilde F^{-4n}.
 \tag{EIR.24}
\]
Its source translation is \(+4nR\). In the good row its source
lies in the identity neighbourhood. In the reduced row the
actual word that removes its order-three component is
\[
 \overline C=C B^{-n}
       =A^{-(m+4)}\mathfrak D^4
       =R_\lambda^{-n(m+4)}\zeta_m^{-4d}
                                  \widetilde F^{-n(m+4)}.
 \tag{EIR.25}
\]
Its source translation is \(+n(m+4)R\), the section in
(EIR.13). The factor \(B^{-n}\) is an actual suspension deck
map. Its integer is retained on the annulus.

On the original root, divide the chosen word by its actual nonzero
central scalar, use the identity formal logarithm, and then add
back the chosen logarithm of that scalar. This gives
\(L_C\) or \(L_{\overline C}\) there. The scalar includes the
entire \(R_\lambda\), finite character and the fixed branch of
the central norm root of four. None is replaced by one.
The resulting directed maps satisfy
\[
 \widehat J_u=\operatorname{Exp}_o(uL_C)J_u,
 \quad \widehat J_{u+1}=A^{-4}\widehat J_u\quad(a=2),
\]
\[
 \widehat J_u=\operatorname{Exp}_o(uL_{\overline C})J_u,
 \quad \widehat J_{u+1}=A^{-(m+4)}\widehat J_u\quad(a=6).
 \tag{EIR.26}
\]
They are closed on the corresponding cap quotients over the
punctured parameter annulus. Their inverse arrows retain the
opposite powers. Transporting these logarithms to native coordinates
uses precisely the domain (EIR.20); the pole (EIR.11) remains in
the naive joint coordinates.

\par

\medskip\noindent\textbf{EIR.6. Literal annular residual and current endpoint of the calculation}

\par\nobreak

Let \(\ell_d=\log d/(2\pi i)\), with its extending chosen logarithm.
The ordinary angular coordinate has the exact change
\[
 s_o=s_e+3\ell_d-4u.
 \tag{EIR.27}
\]
On each sector let \(a_\alpha=n f_\alpha+c_\lambda\delta\)
be the full logarithm of the base-preserving factor in (EIR.23).
All coefficients are transported by the actual point/Cartier map;
(EIR.7) is included when a tensor frame is used. Let \(h_o,h_e\)
be the whole reference cochains in the transported native Lie frame, and
let \(I_e\) be the old-ordinary-to-native-ordinary source/line map,
the inverse of (EIR.2) with its full point/Cartier arrow. Substituting (EIR.27) in the exact sector maps
gives
\[
 \widehat J_u^{-1}\Phi_o=
 \operatorname{Exp}\bigl(3\ell_d a_\alpha
       +u\Pi_C+h_o-h_e\bigr)\Phi_e I_e\quad(a=2),
\]
\[
 \widehat J_u^{-1}\Phi_o=
 \operatorname{Exp}\bigl(3\ell_d a_\alpha
       +u(m a_\alpha+\Pi_{\overline C})+h_o-h_e\bigr)
                       \Phi_e I_e\quad(a=6),
 \tag{EIR.28}
\]
where the full extension-group period columns are
\[
 \Pi_C=-4a_\alpha-L_C,\qquad
 \Pi_{\overline C}=-(m+4)a_\alpha-L_{\overline C}.
 \tag{EIR.29}
\]
In the reduced row \(m a_\alpha=n\operatorname{Log}B+\Lambda_A\),
where \(\Lambda_A\) retains the finite character and the complete
norm-period vector. Thus the extra parameter term includes the
literal \(un\operatorname{Log}B\); it has not been discarded as
an extending root logarithm. Changing a sector branch changes
the displayed period/cochain columns by the corresponding actual
integral suspension and lattice maps.

Equations (EIR.14)-\/-(EIR.29) specify the graph chart, directed cap
word, and annular residual that the remaining attachment must
carry. The following sections add the complementary source/root chart
and the ordinary quotient on its actual domain. The complete
selected-fan attachment still has to carry the full Cartier graph,
including (EIR.22), and the suspension with the actual scalar
\(\lambda_e=e^4\lambda\). The present work assigns no return.

\par

\medskip\noindent\textbf{EIR.7. The complementary graph chart is the actual fourth III root}

\par\nobreak

The other base graph cover is
\[
 z=\xi^4,\qquad e=\xi^m\sigma,
 \qquad X=\xi^{2b_*}\bar X,\quad Y=\xi^{3b_*}\bar Y,
 \quad b_*=r_*+c.
 \tag{EIR.30}
\]
Substitution into (EIR.4), with its full parameter functions,
gives the same equation for both progressions:
\[
 \bar Y^2=\bar X^3+A(e)\xi^{2m}\bar X^2
                    +d(e)\bar X+\sigma^2,
 \qquad \bar G=(0,\sigma).
 \tag{EIR.31}
\]
The identities \(m=3r_*-c\) and \(b_*=r_*+c\) make the
three displayed leading powers exactly equal. The complete
map to the original source is
\[
 x_o=\sigma^2d^{-2}\bar X,\qquad
 y_o=\sigma^3d^{-3}\bar Y,\qquad
 t_o=\sigma^4d^{-3},\qquad
 r=\xi^{2m}\sigma^2A(e)d^{-2}.
 \tag{EIR.32}
\]
This proves the claimed source interpretation by an actual
substitution. At \(\xi=0\), it is the fourth good root of
the original type-III source at \(r=0,t=0\). The joint central
curve is
\[
 E_0:\ \bar Y^2=\bar X^3+\bar X,\qquad
 T=(0,0)=\bar G_0.
 \tag{EIR.33}
\]
The curve at \(\sigma=0\) has equation
\(\bar Y^2=\bar X(\bar X^2-2\xi^{2m}\bar X+1)\).
Its discriminant is nonzero for \(\xi\) in a sufficiently
small disc. Thus the whole elliptic chart is smooth on a
smaller joint bidisc, including its marked zero section.

On the overlap with the old-root graph chart, choose the
displayed root of \(\sigma\) and use
\[
 \eta=\xi\sigma^{1/m},\qquad
 \zeta=\sigma^{-4/m},\qquad z_o=d^{3/m}\sigma^{-4/m}.
 \tag{EIR.34}
\]
For the good row,
\((U,V)=(\sigma^{-2/3}\bar X,\sigma^{-1}\bar Y)\).
For the reduced row,
\[
 U=d^{2/m}\sigma^{-2(l+2)/m}\bar X,\qquad
 V=d^{3/m}\sigma^{-3(l+2)/m}\bar Y.
 \tag{EIR.35}
\]
These equations are the full inverse overlap of
(EIR.15)-\/-(EIR.16), not an identification by fibre type.

The complete \(\mu_4\) action on this cover is
\[
 (\xi,\sigma,\bar X,\bar Y)\longmapsto
 (i\xi,i^{-m}\sigma,i^{-2b_*}\bar X,i^{-3b_*}\bar Y).
 \tag{EIR.36}
\]
It preserves \((e,z,X,Y)\) and all original coordinates.
Its action on the fourth III root is
\(\mathfrak D_{\rm III}^{-m}\), where
\[
 \mathfrak D_{\rm III}:
 (\sigma,\bar X,\bar Y)\longmapsto(i\sigma,-\bar X,i\bar Y)
\]
is its actual root deck. The subgroup of order
\(\gcd(m,4)\) is the generic inertia along \(\xi=0\),
and is retained in this atlas.

The original \(m\)-root deck has the following explicit lift:
\[
 \widetilde{\mathfrak D}:
 (\xi,\sigma,\bar X,\bar Y)\longmapsto
 (\zeta_{4m}\xi,-i\sigma,-\bar X,-i\bar Y).
 \tag{EIR.37}
\]
It preserves \(e\), sends \(z\) to \(\zeta_mz\), and sends
the native \((X,Y)\) to their original root weights
\(\zeta_m^{-2k}X,\zeta_m^{-3k}Y\). Its \(m\)-th power
is precisely (EIR.36). Therefore both the graph cover and the
actual finite-cap root action are specified at this corner.

\par

\medskip\noindent\textbf{EIR.8. Full selected Hom factor and the finite division cocycle}

\par\nobreak

In (EIR.32) the selected factor is
\[
 t_o-t_{c,o}=\frac{\sigma^4}{d^3}H,\qquad
 H=1-\xi^{4m}(1-e^2).
\]
The actual original Miller functions have weight
\((\sigma/d)^{a-1}\). Their six-factor source/target
composition consequently gives the exact line arrow
\[
 B_{{\rm sel},o}=H\sigma^{-2}\bar B_{a,b}.
 \tag{EIR.38}
\]
Here \(\bar B_{a,b}\) is the full point-line coefficient
on (EIR.31), formed by the same actual Miller source and
target factors. Its generic central orders can be evaluated
without replacing the point line by a tensor power. On a
curve \(y^2=x^3+a_2x^2+d x+\sigma^2\), the tangent at
\((0,\sigma)\) has slope \(d/(2\sigma)\), and its
second point has \(x=d^2/(4\sigma^2)-a_2\). Thus
\[
 \bar f_2/\sigma\longrightarrow 2\bar X/d,
 \qquad \bar f_3\longrightarrow\bar X,
 \qquad \bar f_6/\sigma\longrightarrow6\bar X^3/d.
 \tag{EIR.39}
\]
The second limit is the product of the two full chord-over-
vertical functions; the third retains the actual tangent
at \(3G\), whose leading slope is \(d/(6\sigma)\).
The central translation by \(-G\) is
\[
 (\bar X,\bar Y)\longmapsto
 (d/\bar X,-d\bar Y/\bar X^2).
\]
Substitution into the actual M6 source/target ratio and
into all four M2 orbit factors yields
\[
 \bar B_{6,1}/\sigma^6\longrightarrow64/d^3,
 \qquad \bar B_{2,3}/\sigma^6\longrightarrow64/d^3.
 \tag{EIR.40}
\]
In (EIR.31), \(d|_{\sigma=0}=1\), so the coefficient is
64. The whole Cartier/normal extension of (EIR.38),
including its charts at both sections and the finite
action (EIR.37), is being written by fourfold. The full
factor carried into that calculation is
\(H\sigma^4(\bar B_{a,b}/\sigma^6)\); the power of
\(\sigma\) has not been absorbed into a unit.

There is also a concrete finite division constraint.
The central source map of (EIR.37) is
\(\mathfrak D_{\rm III}^{-1}\) on (EIR.33).
Its fixed points are exactly \(O,T\): a finite fixed
point has \(-\bar X=\bar X\) and \(-i\bar Y=\bar Y\),
hence \(\bar X=\bar Y=0\). For even \(m\), neither fixed
point has \(mR=bT=T\), because \(b=1\) or 3 is odd.
Thus an ordinary equivariant division section cannot
extend through this point. This calculation identifies
the nontrivial finite torsion cocycle that the actual
root-groupoid continuation must carry; it does not
exclude a construction with that cocycle.

For the smallest reduced order \(m=2\), any central
half \(R_0\) of \(T\) satisfies
\[
 \mathfrak D_{\rm III}^{-2}R_0-R_0
                  =-2R_0=T.
 \tag{EIR.41}
\]
For example, the four halves have \(x=1,-1\) and
the corresponding \(y=\pm\sqrt2,\pm i\sqrt2\).
On a local \(m=2\) division chart, let \(T(\xi,\sigma)\)
be the actual two-torsion section with central value
\(T\); it is the analytic cubic root near \(x=0\) with
\(y=0\). The curve coefficients are functions of
\(\xi^4,\sigma^2\) in this case. The unique local
division branch with value \(R_0\) then obeys
\[
 R(\xi,-\sigma)=-R(\xi,\sigma)+T(\xi,\sigma).
 \tag{EIR.42}
\]
Both sides have double \(-G(\xi,\sigma)\), and they
have the same central value; the actual local inverse
of \([2]\) proves the equality. This torsion section
agrees with \(G\) only on the central fibre. Replacing
it by \(G\) off that fibre would discard an actual
source translation and is not done here.

\par

\medskip\noindent\textbf{EIR.9. The actual cap domain and the ordinary quotient of the corner}

\par\nobreak

The source map (EIR.32) locates the action required at each point:
\[
 |q_{\infty,o}|=|d|^3/|\sigma|^4.
\]
Hence an original finite-cap domain \(|q_{\infty,o}|<\delta\)
has exactly \(|\sigma|>|d|^{3/4}\delta^{-1/4}\).
It excludes \(\sigma=0\). The original K96 cap continues on
that actual domain, carrying its full norm, character and reference
cochain. The finite-division obstruction in EIR.8 concerns an attempted
extension of that cap action through the larger corner; it is not an
obstruction assigned to the ordinary-source piece required there.
The latter piece carries the actual suspension \(B\).

The subgroup generated by \(\widetilde{\mathfrak D}^{\,4}\)
in (EIR.37) is \(\mu_m\). It acts by
\(\xi\mapsto\zeta_m\xi\) and fixes
\(\sigma,\bar X,\bar Y\). The actual original point-line
pullback is the identity under this subgroup, because the original
source point and its canonical Cartier section are both unchanged.
The full invariant coordinate map is therefore
\[
 \chi=\xi^m,\qquad
 e=\chi\sigma,\qquad q_{\infty,e}=\chi^4,
 \qquad x_e=\chi^{-2}\bar X,\quad
 y_e=\chi^{-3}\bar Y.
 \tag{EIR.43}
\]
Its equation and remaining action are
\[
 \bar Y^2=\bar X^3+A(e)\chi^2\bar X^2
                    +d(e)\bar X+\sigma^2,
 \quad \bar G=(0,\sigma),
\]
\[
 (\chi,\sigma,\bar X,\bar Y)\longmapsto
           (i\chi,-i\sigma,-\bar X,-i\bar Y).
 \tag{EIR.44}
\]
This is a literal quotient: the invariant algebra is generated by
\(\xi^m,\sigma,\bar X,\bar Y\), with the displayed equation.
The original \(m\)-division section is not used on this ordinary
chart. Every \(m\)-cap character and its full annular cochain remain
on the actual overlap with its original cap domain.

The base relation and the full ordinary suspension become
\[
 \chi^4t_o=e^4/d^3,\qquad
 B=\lambda H\sigma^{-2}\bar B_{a,b},\qquad
 H=1-\chi^4(1-e^2).
 \tag{EIR.45}
\]
Thus the selected zero is exactly \(H=0\); it is retained along
with the power \(\sigma^4\) in (EIR.40). The source at the
joint origin is a smooth elliptic curve before the remaining
\(\mu_4\) quotient. The fixed normal/source line arrows for that
quotient, together with the actual selected-node toric charts,
are the current ordinary attachment calculation. The complete
weighted root atlas has not been replaced by a native cap on
a larger domain.

\par

\medskip\noindent\textbf{EIR.10. Actual point-line extension and the ordinary chain}

\par\nobreak

Here all coordinates are those of the ordinary quotient (EIR.43),
and \(H\ne0\). Fourfold\textquotesingle s full Cartier calculation in EI.7 gives
the following literal section maps. On the smooth joint elliptic
family write \(T_b=t_{-b\bar G}\). The actual Miller functions satisfy
\[
 \operatorname{div}\bar f_2
   =2[\bar G]-[2\bar G]-[O]+[\sigma=0],
\qquad
 \operatorname{div}\bar f_6
   =6[\bar G]-[6\bar G]-5[O]+[\sigma=0].
 \tag{EIR.46}
\]
Their complete chronological products, with the source and target
point frames carried by \(T_b\), therefore have
\[
 \operatorname{div}\bar B_{a,b}
 =[a\bar G]+[b\bar G]-[(a+b)\bar G]-[O]+6[\sigma=0]
 \quad ((a,b)=(6,1),(2,3)).
 \tag{EIR.47}
\]
This keeps the divisors where the four horizontal sections meet.
For example, the \(M_6\) product is
\(\bar f_2^6\,\bar f_6(T_1p)/\bar f_6(p)\); the \(M_2\)
product consists of
\(\bar f_2(p)\bar f_2(T_3p)
  \bar f_2(T_1p)^2\bar f_2(T_2p)^2\).
Substitution of (EIR.46) gives (EIR.47) term by term.

Let \(\sigma_a\) be the canonical meromorphic section of the
actual point line \(M_a\). In a neighbourhood of \(O\) put
\[
 r_a(p)=\frac{\tau(p)-\tau(a\bar G)}{\tau(p)},\qquad
 e_a=\sigma_a/r_a .
 \tag{EIR.48}
\]
The coordinate \(\tau(a\bar G)\) extends on the whole smaller
joint base disc. For \(a=2\) this follows from the exact expression
\[
 \tau(2\bar G)=\frac{2\sigma S}{dS+2\sigma^4},
 \quad
 S=\frac{d^2}{4}-A(e)\chi^2\sigma^2
   =\frac{(1-e^2)(1+3e^2)}4,
 \tag{EIR.49}
\]
whose denominator is a unit. For \(a=6\), the group-law morphism
\([3]\) on the same smooth elliptic family carries this small
section to \(6\bar G\); its identity chart is holomorphic and has
derivative \(3\). Outside the two small horizontal section
neighbourhoods use their canonical frame. The frames in (EIR.48)
and their translates by \(T_b\) cover every coalescing horizontal
section.

In these genuine nonzero line frames the normalized suspension is
the full map with coefficient
\[
 \mathcal B_{a,b}(p)
   =\sigma^{-6}\bar B_{a,b}(p)
       \frac{r_a(T_b p)}{r_a(p)} .
 \tag{EIR.50}
\]
The divisor of the last ratio is the negative of the four
horizontal terms in (EIR.47). Hence both this coefficient and its
reciprocal have zero divisor in every displayed smooth chart.
Normality extends them across the intersections of those divisors.
This proves the full invertible Hom map, with central value \(64\)
in the carried frames, instead of inferring extension from its
generic value. The actual deck carries \(\tau\) and
\(\tau(a\bar G)\) with the same character, so it preserves
(EIR.48) and the point-line coefficient. It also preserves
\(\mathcal B_{a,b}\).

The ordinary line charts are now literal:
\[
 X_jY_j=\sigma,\qquad V=\sigma^jX_j,\qquad j\in\mathbb Z.
 \tag{EIR.51}
\]
Here \(V\) means the coefficient in the genuine point frame just
specified; the full changes of point frame remain in the
transition functions. On the overlap of consecutive charts,
\[
 X_{j+1}=Y_j^{-1},\qquad Y_{j+1}=\sigma Y_j.
 \tag{EIR.52}
\]
The full suspension and remaining finite action are
\[
 \begin{aligned}
 B:\ (p,j,X_j,Y_j)&\longmapsto
 \left(T_bp,j+4,\lambda H\mathcal B_{a,b}(p)X_j,
       (\lambda H\mathcal B_{a,b}(p))^{-1}Y_j\right),\\
 \gamma:\ (\chi,\sigma,p,X_j,Y_j)&\longmapsto
 \left(i\chi,-i\sigma,\mathfrak D_{\rm III}^{-1}p,
                  i^jX_j,(-i)i^{-j}Y_j\right).
 \end{aligned}
 \tag{EIR.53}
\]
They preserve the equation in (EIR.51). Their commutation is
literal: the source maps commute, the normalized Hom coefficient
is deck invariant, and \(i^{j+4}=i^j\). The action of \(B\) moves
the integer index by four even at the central nodes. Thus these
charts give the actual ordinary suspension quotient on the
\(H\)-unit patch.

On the overlap with the original finite cap, \(\sigma\) is
nonzero by EIR.9. Equations (EIR.51)-\/-(EIR.52) are then just
the original point line with its displayed coordinate \(V\).
The cap comparison is exactly (EIR.28), including the signed
period columns, scalar powers, finite character and the
additional \(B^{-n}\) reindex in the reduced row. No extra
choice of annular integer is inserted by (EIR.51).

The locus \(H=0\) is excluded from the inversion of the coefficient
in (EIR.53). It retains the actual selected-node periodic fan.
These formulas prove the ordinary \(H\)-unit attachment and its
overlap with the original cap; they do not yet give a compact
selected-fan filling or a numerical return.

\par

\medskip\noindent\textbf{EIR.11. Finite-native chain overlap and the full old-cap groupoid square}

\par\nobreak

On the \(\chi\)-unit overlap the complete ordinary source map is
\[
 (x_e,y_e,t_e)=(\chi^{-2}\bar X,\chi^{-3}\bar Y,\chi^{-4}),
 \qquad e=\chi\sigma,\qquad q_e=\chi^{-4}H .
 \tag{EIR.54}
\]
The canonical point section is carried literally. The Miller functions
and the complete six-factor Hom coefficient obey
\[
 f_{a,e}=\chi^{1-a}\bar f_a,\qquad
 B_{a,b,e}=\chi^{-6}\bar B_{a,b}.
\]
Consequently the exact selected Hom arrow is
\[
 \widehat B_e
   =e^{-6}(q_e/t_e^3)B_{a,b,e}
   =\chi^{-4}H\,\sigma^{-6}\bar B_{a,b},
 \qquad
 \lambda e^4\widehat B_e
   =\lambda\sigma^4H\,\sigma^{-6}\bar B_{a,b}.
 \tag{EIR.55}
\]
All equations are identities of the actual source and target line
frames. The nonzero point frames also agree:
\(\tau_e=\chi\bar\tau\) and
\(\tau_e(aG)=\chi\bar\tau(aG)\), so their Cartier ratios agree
identically. For native line charts
\[
 A_jB_j=e,\qquad V=e^jA_j,
\]
the ordinary overlap is therefore
\[
 X_j=\chi^jA_j,\qquad Y_j=\chi^{-j-1}B_j.
 \tag{EIR.56}
\]
It preserves both \(X_jY_j=\sigma\) and the physical coefficient
\(V\). It gives exactly the finite phases in (EIR.53). The
finite-native \(I_2\) source resolution and its full translation
maps in EI.15-\/-17 are carried on this overlap; no extra line
integer is inserted.

Here is the complete old-cap square on its actual large-\(\sigma\)
overlap. Write
\[
 v_\sigma=\frac{\log\sigma}{2\pi i},\qquad
 s_o=\frac{\log q_{\infty,o}}{2\pi i}
     =3\ell_d-4v_\sigma .
 \tag{EIR.57}
\]
The graph generator \(\gamma\) carries
\((\chi,\sigma)\) to \((i\chi,-i\sigma)\), with \(e\) fixed.
On these logarithmic sectors it sends
\(v_\sigma\) to \(v_\sigma-\tfrac14\), hence \(s_o\) to
\(s_o+1\). The original ordinary source point and its canonical
point-line coefficient are unchanged by this graph map.

Keep exactly the K96 EX.17-\/-19 cochains on that original cap.
In the original root Lie frame they are
\[
 \begin{aligned}
 \Phi_\alpha&=
 \operatorname{Exp}\bigl(ns_\alpha f_\alpha
              +s_\alpha c_\lambda\delta+h_\alpha\bigr)I_{z_\alpha},\\
 f_\beta&=D_*^{\,k}f_\alpha+\Pi_{\alpha\beta},\\
 h_\beta-D_*^{\,k}h_\alpha
   &=-ns_\alpha\Pi_{\alpha\beta}
          +j_{\alpha\beta}\operatorname{Log}_\beta B
          +\Lambda_{\alpha\beta},\qquad
 s_\beta=s_\alpha+k .
 \end{aligned}
 \tag{EIR.58}
\]
The functions, finite characters and all integers in this display
are the prescribed transported reference data; \(n=\pm1\) is
retained. The literal squares are
\[
 \Phi_\beta=A^kB^{j_{\alpha\beta}}\Phi_\alpha,\qquad
 \Phi_\alpha B=B\Phi_\alpha,\qquad A^m=B^n .
 \tag{EIR.59}
\]
Thus the root representative of the quarter-turn graph action is
\(AB^{j_\gamma}\), with its actual reference integer, and the
four-step representative is
\[
 A^4 B^{\,J_\gamma},\qquad
 J_\gamma=\sum_{\text{four traversed sectors}}j_{\alpha\beta}.
 \tag{EIR.60}
\]
Every term in this sum comes from the displayed sector continuation.
The full \(D\)-transport of \(\Pi,\Lambda\) remains in (EIR.58).
On the cap quotient the \(A,B\) deck words are identified, so
this is precisely the finite fourth-root graph action there.
On the root groupoid the words (EIR.59)-\/-(EIR.60) remain literal;
neither \(j_\gamma=0\) nor \(J_\gamma=0\) is asserted.
Combining this square with (EIR.28) retains the earlier reduced
row \(B^{-n}\) reindex and its full parameter-period column.

The K96 source/root and Hom maps extend over the original
parameter disc \(r=e^2(3e^2-2)/d^2\) on this old cap domain.
The moving selected height is a unit on its fixed infinity
annulus; its entire value is retained in \(B_{\rm sel}\) and
in (EIR.58). Consequently these formulas glue that unchanged
cap to the ordinary graph on the displayed overlap. They
require no division section on the ordinary chart at
\(\sigma=0\). The selected toric/source graph assembly still
retains its own finite strata and all boundary components.

